\documentclass[12pt]{exam}
\usepackage[utf8]{inputenc}
\usepackage{amsmath, amssymb}
\usepackage{graphicx}
\usepackage[export]{adjustbox}
\usepackage[margin=0.5in, landscape]{geometry}
\usepackage{multicol}
\usepackage{listings}
\usepackage{xcolor}
\pagecolor{black}
\color{white}
\pagestyle{empty}
\cfoot{}

\lstset{
    basicstyle=\ttfamily\small\color{white},
    backgroundcolor=\color[rgb]{0.12,0.12,0.12},
    frame=single,
    rulecolor=\color{gray},
    breaklines=true,
    tabsize=4
}

\begin{document}
\begin{center}
\LARGE \textbf{GATE CSE: 1. Discrete Maths - 1. Mathematical Logic}
\end{center}
\vspace{0.5cm}
\begin{questions}
\begin{multicols*}{2}

\question {\small \textbf{[GATE 1995 | 2 Mark]}} \\[0.2cm]
{\Large If the proposition $\neg p \Rightarrow q$ is true, then the truth value of the proposition $\neg p \vee (p \Rightarrow q)$, where $\neg$ is negation, $\vee$ is inclusive or and $\Rightarrow$ is implication, is}
\begin{choices}
  \choice True
  \choice Multiple valued
  \choice False
  \choice Cannot be determined
\end{choices}

\end{multicols*}
\vspace{0.6cm}
\newpage
\begin{multicols*}{2}

\question {\small \textbf{[GATE 1996 | 2 Mark]}} \\[0.2cm]
{\Large Which one of the following is false? Read $\wedge$ as AND, $\vee$ as OR, $\sim$ as NOT, $\rightarrow$ as one way implication and $\leftrightarrow$ as two way implication.}
\begin{choices}
  \choice $((x \rightarrow y) \wedge x) \rightarrow y$
  \choice $((\sim x \rightarrow y) \wedge (\sim x \rightarrow \sim y)) \rightarrow x$
  \choice $(x \rightarrow (x \vee y))$
  \choice $((x \vee y) \leftrightarrow (\sim x \rightarrow \sim y))$
\end{choices}

\end{multicols*}
\vspace{0.6cm}
\newpage
\begin{multicols*}{2}

\question {\small \textbf{[GATE 2000 | 2 Mark]}} \\[0.2cm]
{\Large Let $a, b, c, d$ be propositions. Assume that the equivalence $a \leftrightarrow (b \vee \neg b)$ and $b \leftrightarrow c$ hold. Then the truth-value of the formula $(a \wedge b) \rightarrow (a \wedge c) \vee d$ is always}
\begin{choices}
  \choice True
  \choice False
  \choice Same as the truth-value of $b$
  \choice Same as the truth-value of $d$
\end{choices}

\end{multicols*}
\vspace{0.6cm}
\newpage
\begin{multicols*}{2}

\question {\small \textbf{[GATE 2001 | 1 Mark]}} \\[0.2cm]
{\Large What is the converse of the following assertion? I stay only if you go}
\begin{choices}
  \choice I stay if you go
  \choice If I stay then you go
  \choice If you do not go then I do not stay
  \choice If I do not stay then you go
\end{choices}

\end{multicols*}
\vspace{0.6cm}
\newpage
\begin{multicols*}{2}

\question {\small \textbf{[GATE 2001 | 1 Mark]}} \\[0.2cm]
{\Large Consider two well-formed formulas in propositional logic: $F_1 : P \Rightarrow \neg P$ $F_2 : (P \Rightarrow \neg P) \vee (\neg P \Rightarrow P)$ Which of the following statements is correct?}
\begin{choices}
  \choice $F_1$ is satisfiable, $F_2$ is valid
  \choice $F_1$ is unsatisfiable, $F_2$ is satisfiable
  \choice $F_1$ is unsatisfiable, $F_2$ is valid
  \choice $F_1$ and $F_2$ are both satisfiable
\end{choices}

\end{multicols*}
\vspace{0.6cm}
\newpage
\begin{multicols*}{2}

\question {\small \textbf{[GATE 2002 | 1 Mark]}} \\[0.2cm]
{\Large "If X then Y unless Z" is represented by which of the following formulas in propositional logic? ("$\neg$" is negation, "$\wedge$" is conjunction, and "$\rightarrow$" is implication)}
\begin{choices}
  \choice $(X \wedge \neg Z) \rightarrow Y$
  \choice $(X \wedge Y) \rightarrow \neg Z$
  \choice $X \rightarrow (Y \wedge \neg Z)$
  \choice $(X \rightarrow Y) \wedge \neg Z$
\end{choices}

\end{multicols*}
\vspace{0.6cm}
\newpage
\begin{multicols*}{2}

\question {\small \textbf{[GATE 2003 | 2 Mark]}} \\[0.2cm]
{\Large Which of the following is a valid first order formula? (Here $\alpha$ and $\beta$ are first order formulae with x as their only free variable)}
\begin{choices}
  \choice $((\forall x) [\alpha] \Rightarrow (\forall x) [\beta]) \Rightarrow (\forall x) [\alpha \Rightarrow \beta]$
  \choice $(\forall x) [\alpha] \Rightarrow (\exists x) [\alpha \wedge \beta]$
  \choice $(\forall x) [\alpha \vee \beta] \Rightarrow (\exists x) [\alpha] \Rightarrow (\forall x) [\alpha]$
  \choice $(\forall x) [\alpha \Rightarrow \beta] \Rightarrow ((\forall x) [\alpha] \Rightarrow (\forall x) [\beta])$
\end{choices}

\end{multicols*}
\vspace{0.6cm}
\newpage
\begin{multicols*}{2}

\question {\small \textbf{[GATE 2003 | 2 Mark]}} \\[0.2cm]
{\Large Consider the following formula $\alpha$ and its two interpretations $I_1$ and $I_2$.\newline
$\alpha : (\forall x) [P_x \Leftrightarrow (\forall y) [Q_{xy} \Leftrightarrow \neg Q_{yy}]] \Rightarrow (\forall x) [\neg P_x]$\newline
 $I_1 :$ Domain: the set of natural numbers \newline
$P_x =$ 'x is a prime number' \newline
$Q_{xy} =$ 'y divides x' \newline
$I_2 :$ Same as $I_1$ except that $P_x =$ 'x is a composite number.' \newline
Which of the following statements is true?}
\begin{choices}
  \choice $I_1$ satisfies $\alpha$, $I_2$ does not
  \choice $I_2$ satisfies $\alpha$, $I_1$ does not
  \choice Neither $I_2$ nor $I_1$ satisfies $\alpha$
  \choice Both $I_1$ and $I_2$ satisfy $\alpha$
\end{choices}

\end{multicols*}
\vspace{0.6cm}
\newpage
\begin{multicols*}{2}

\question {\small \textbf{[GATE 2003 | 2 Mark]}} \\[0.2cm]
{\Large The following resolution rule is used in logic programming: Derive clause $(P \vee Q)$ from clauses $(P \vee R), (Q \vee \neg R)$ Which of the following statemnets related to this rule is FALSE?}
\begin{choices}
  \choice $(P \vee R) \wedge (Q \vee \neg R) \Rightarrow (P \vee Q)$ is logically valid
  \choice $(P \vee Q) \Rightarrow (P \vee R) \wedge (Q \vee \neg R)$ is logically valid
  \choice $(P \vee Q)$ is satisfiable if and only if $(P \vee R) \wedge (Q \vee \neg R)$ is satisfiable
  \choice $(P \vee Q) \Rightarrow FALSE$ if and only if both P and Q are unsatisfiable
\end{choices}

\end{multicols*}
\vspace{0.6cm}
\newpage
\begin{multicols*}{2}

\question {\small \textbf{[GATE 2004 | 1 Mark]}} \\[0.2cm]
{\Large Identify the correct translation into logical notation of the following assertion. Some boys in the class are taller than all the girls\newline
$\textbf{Note:}$ Taller (x, y) is true if x is taller than y.}
\begin{choices}
  \choice $(\exists x) (\text{boy}(x) \to (\forall y) (\text{girl}(y) \wedge \text{taller} (x, y)))$
  \choice $(\exists x) (\text{boy}(x) \wedge (\forall y) (\text{girl}(y) \wedge \text{taller} (x, y)))$
  \choice $(\exists x) (\text{boy}(x) \to (\forall y) (\text{girl}(y) \to \text{taller} (x, y)))$
  \choice $(\exists x) (\text{boy}(x) \wedge (\forall y) (\text{girl}(y) \to \text{taller} (x, y)))$
\end{choices}

\end{multicols*}
\vspace{0.6cm}
\newpage
\begin{multicols*}{2}

\question {\small \textbf{[GATE 2004 | 1 Mark]}} \\[0.2cm]
{\Large Let $a(x, y)$, $b(x, y)$, and $c(x, y)$ be three statements with variables $x$ and $y$ chosen from some universe. Consider the following statement:\newline
 $(\exists x)(\forall y)[(a(x, y) \wedge b(x, y)) \wedge \neg c(x, y)]$\newline
Which one of the following is its equivalent?}
\begin{choices}
  \choice $(\forall x)(\exists y)[(a(x, y) \vee b(x, y)) \to c(x, y)]$
  \choice $(\exists x)(\forall y)[(a(x, y) \vee b(x, y)) \wedge \neg c(x, y)]$
  \choice $\neg (\forall x)(\exists y)[(a(x, y) \wedge b(x, y)) \to c(x, y)]$
  \choice $\neg (\forall x)(\exists y)[(a(x, y) \vee b(x, y)) \to c(x, y)]$
\end{choices}

\end{multicols*}
\vspace{0.6cm}
\newpage
\begin{multicols*}{2}

\question {\small \textbf{[GATE 2004 | 2 Mark]}} \\[0.2cm]
{\Large Let $p, q, r$ and $s$ be four primitive statements. Consider the following arguments:\newline
P: $[(\neg p \vee q) \wedge (r \to s) \wedge (p \vee r)] \to (\neg s \to q)$ \newline
Q: $[(\neg p \wedge q) \wedge [q \to (p \to r)]] \to \neg r$ \newline
R: $[[(q \wedge r) \to p] \wedge (\neg q \vee p)] \to r$ S: $[p \wedge (p \to r) \wedge (q \vee \neg r)] \to q$ \newline
Which of the above arguments are valid?}
\begin{choices}
  \choice P and Q only
  \choice P and R only
  \choice P and S only
  \choice P, Q, R and S
\end{choices}

\end{multicols*}
\vspace{0.6cm}
\newpage
\begin{multicols*}{2}

\question {\small \textbf{[GATE 2004 | 2 Mark]}} \\[0.2cm]
{\Large The following propositional statement is \newline
$(P \to (Q \vee R)) \to ((P \wedge Q) \to R)$}
\begin{choices}
  \choice Satisfiable but not valid
  \choice Valid
  \choice A contraditiction
  \choice None of the above
\end{choices}

\end{multicols*}
\vspace{0.6cm}
\newpage
\begin{multicols*}{2}

\question {\small \textbf{[GATE 2005 | 2 Mark]}} \\[0.2cm]
{\Large Let P, Q and R be three atomic propositional assertions. Let X denote $(P \vee Q) \to R$ and Y denote $(P \to R) \vee (Q \to R)$. Which one of the following is a tautology?}
\begin{choices}
  \choice $X \equiv Y$
  \choice $X \to Y$
  \choice $Y \to X$
  \choice $\neg Y \to X$
\end{choices}

\end{multicols*}
\vspace{0.6cm}
\newpage
\begin{multicols*}{2}

\question {\small \textbf{[GATE 2005 | 2 Mark]}} \\[0.2cm]
{\Large What is the first order predicate calculus statement equivalent to the following? Every teacher is liked by some student}
\begin{choices}
  \choice $\forall(x) [\text{teacher} (x) \to \exists(y) [\text{student} (y) \to \text{likes} (y, x)]]$
  \choice $\forall(x) [\text{teacher} (x) \to \exists(y) [\text{student} (y) \wedge \text{likes} (y, x)]]$
  \choice $\exists(y) \forall(x) [\text{teacher} (x) \to [\text{student} (y) \wedge \text{likes} (y, x)]]$
  \choice $\forall(x) [\text{teacher} (x) \wedge \exists(y) [\text{student} (y) \to \text{likes} (y, x)]]$
\end{choices}

\end{multicols*}
\vspace{0.6cm}
\newpage
\begin{multicols*}{2}

\question {\small \textbf{[GATE 2005 | 2 Mark]}} \\[0.2cm]
{\Large Let $P(x)$ and $Q(x)$ be arbitrary predicates. Which of the following statements is always TRUE?}
\begin{choices}
  \choice $(\forall x(P(x) \vee Q(x))) \Rightarrow ((\forall x P(x)) \vee (\forall x Q(x)))$
  \choice $(\forall x(P(x) \Rightarrow Q(x))) \Rightarrow ((\forall x P(x)) \Rightarrow (\forall x Q(x)))$
  \choice $(\forall x(P(x) \Rightarrow (\forall x Q(x)))) \Rightarrow (\forall x(P(x) \Rightarrow Q(x)))$
  \choice $((\forall x(P(x)) \Leftrightarrow (\forall x Q(x))) \Rightarrow (\forall x(P(x) \Leftrightarrow Q(x)))$
\end{choices}

\end{multicols*}
\vspace{0.6cm}
\newpage
\begin{multicols*}{2}

\question {\small \textbf{[GATE 2006 | 2 Mark]}} \\[0.2cm]
{\Large Consider the following first order logic formula in which $R$ is a binary relation symbol.\newline
$\forall x \forall y (R(x, y) \Rightarrow R(y, x))$ The formula is}
\begin{choices}
  \choice Satisfiable and valid
  \choice Satisfiable and so is its negation
  \choice Unsatisfiable but its negation is valid
  \choice Satisfiable but its negation is unsatisfiable
\end{choices}

\end{multicols*}
\vspace{0.6cm}
\newpage
\begin{multicols*}{2}

\question {\small \textbf{[GATE 2006 | 2 Mark]}} \\[0.2cm]
{\Large Which one of the first order predicate calculus statements given below correctly expresses the following English statement?\newline
Tigers and lions attack if they are hungry or threatened.}
\begin{choices}
  \choice $\forall x [(\text{tiger} (x) \wedge \text{lion} (x)) \to \{(\text{hungry} (x) \vee \text{threatened} (x)) \to \text{attacks} (x)\}]$
  \choice $\forall x [(\text{tiger} (x) \vee \text{lion} (x)) \to \{(\text{hungry} (x) \vee \text{threatened} (x)) \wedge \text{attacks} (x)\}]$
  \choice $\forall x [(\text{tiger} (x) \vee \text{lion} (x)) \to \{(\text{attacks} (x) \to (\text{hungry} (x) \vee \text{Threatened} (x))\}]$
  \choice $\forall x [(\text{tiger} (x) \vee \text{lion} (x)) \to \{(\text{hungry} (x) \vee \text{threatened} (x)) \to \text{attacks} (x)\}]$
\end{choices}

\end{multicols*}
\vspace{0.6cm}
\newpage
\begin{multicols*}{2}

\question {\small \textbf{[GATE 2006 | 2 Mark]}} \\[0.2cm]
{\Large Consider the following propositional statements: \newline
$P_1 : ((A \wedge B) \to C) \equiv ((A \to C) \wedge (B \to C))$\newline
$P_2 : ((A \vee B) \to C) \equiv ((A \to C) \vee (B \to C))$ \newline
Which one of the following is true?}
\begin{choices}
  \choice $P_1$ is a tautology, but not $P_2$
  \choice $P_2$ is a tautology, but not $P_1$
  \choice $P_1$ and $P_2$ are both tautologies
  \choice Both $P_1$ and $P_2$ are not tautologies
\end{choices}

\end{multicols*}
\vspace{0.6cm}
\newpage
\begin{multicols*}{2}

\question {\small \textbf{[GATE 2006 | 2 Mark]}} \\[0.2cm]
{\Large A logical binary relation $\odot$, is defined as follows: 
\begin{center}
\includegraphics[max width=0.9\linewidth, keepaspectratio]{"E:/gate-pyq-app/public/1. Discrete Maths/1. Mathematical Logic/20.png"}
\end{center}
 Let $\sim$ be the unary negation (NOT) operator, with higher precedence, than $\odot$. Which one of the following is equivalent to $A \wedge B$?}
\begin{choices}
  \choice $(\sim A \odot B)$
  \choice $\sim (A \odot \sim B)$
  \choice $\sim (\sim A \odot \sim B)$
  \choice $\sim (\sim A \odot B)$
\end{choices}

\end{multicols*}
\vspace{0.6cm}
\newpage
\begin{multicols*}{2}

\question {\small \textbf{[GATE 2007 | 2 Mark]}} \\[0.2cm]
{\Large Let Graph ($x$) be a predicate which denotes that $x$ is a graph. Let Connected ($x$) be a predicate which denotes that $x$ is connected. Which of the following first order logic sentences DOES NOT represent the statement; "Not every graph is connected"?}
\begin{choices}
  \choice $\neg \forall x (\text{Graph}(x) \Rightarrow \text{Connected}(x))$
  \choice $\exists x (\text{Graph}(x) \wedge \neg \text{Connected}(x))$
  \choice $\neg \forall x (\neg \text{Graph}(x) \vee \text{Connected}(x))$
  \choice $\forall x (\text{Graph}(x) \Rightarrow \neg \text{Connected}(x))$
\end{choices}

\end{multicols*}
\vspace{0.6cm}
\newpage
\begin{multicols*}{2}

\question {\small \textbf{[GATE 2007 | 2 Mark]}} \\[0.2cm]
{\Large Which of the following is TRUE about formulae in Conjunctive Normal Form?}
\begin{choices}
  \choice For any formula, there is a truth assignment for which at least half the clauses evaluate to true.
  \choice For any formula, there is a truth assignment for which all the clauses evaluate to true.
  \choice There is a formula such that for each truth assignment at most one-fourth of the clauses evaluate to true.
  \choice None of the above
\end{choices}

\end{multicols*}
\vspace{0.6cm}
\newpage
\begin{multicols*}{2}

\question {\small \textbf{[GATE 2007 | 2 Mark]}} \\[0.2cm]
{\Large Which one of these first-order logic formulae is valid?}
\begin{choices}
  \choice $\forall x(P(x) \Rightarrow Q(x)) \Rightarrow ((\forall xP(x)) \Rightarrow (\forall xQ(x)))$
  \choice $\exists x(P(x) \vee Q(x)) \Rightarrow ((\exists xP(x)) \Rightarrow (\exists xQ(x)))$
  \choice $\exists x(P(x) \wedge Q(x)) \Leftrightarrow ((\exists xP(x)) \wedge (\exists xQ(x)))$
  \choice $\forall x \exists y P(x, y) \Rightarrow \exists y \forall x P(x, y)$
\end{choices}

\end{multicols*}
\vspace{0.6cm}
\newpage
\begin{multicols*}{2}

\question {\small \textbf{[GATE 2008 | 1 Mark]}} \\[0.2cm]
{\Large Let $fsa$ and $pda$ be two predicates such that $fsa(x)$ means $x$ is a finite state automaton, and $pda(y)$ means, that $y$ is a pushdown automaton. Let $equivalent$ be another predicate such that $equivalent (a, b)$ means $a$ and $b$ are equivalent. Which of the following first order logic statement represents the following:\newline
Each finite state automaton has an equivalent pushdown automaton.}
\begin{choices}
  \choice $\forall x (fsa(x) \Rightarrow \exists y (pda(y) \wedge equivalent (x, y)))$
  \choice $\sim \forall y (\exists x fsa(x) \Rightarrow pda(y) \wedge equivalent (x, y))$
  \choice $\forall x \exists y (fsa(x) \wedge pda(y) \wedge equivalent (x, y))$
  \choice $\forall x \exists y (fsa(y) \wedge pda(x) \wedge equivalent (x, y))$
\end{choices}

\end{multicols*}
\vspace{0.6cm}
\newpage
\begin{multicols*}{2}

\question {\small \textbf{[GATE 2008 | 2 Mark]}} \\[0.2cm]
{\Large Which of the following first order formulae is logically valid? Here $\alpha(x)$ is a first order formula with $x$ as a free variable, and $\beta$ is a first order formula with no free variable.}
\begin{choices}
  \choice $[\beta \rightarrow (\exists x, \alpha(x))] \rightarrow [\forall x, \beta \rightarrow \alpha(x)]$
  \choice $[\exists x, \beta \rightarrow \alpha(x)] \rightarrow [\beta \rightarrow (\forall x, \alpha(x))]$
  \choice $[(\exists x, \alpha(x)) \rightarrow \beta] \rightarrow [\forall x, \alpha(x) \rightarrow \beta]$
  \choice $[(\forall x, \alpha(x)) \rightarrow \beta] \rightarrow [\forall x, \alpha(x) \rightarrow \beta]$
\end{choices}

\end{multicols*}
\vspace{0.6cm}
\newpage
\begin{multicols*}{2}

\question {\small \textbf{[GATE 2008 | 2 Mark]}} \\[0.2cm]
{\Large Which of the following is the negation of $[\forall x, \alpha \rightarrow (\exists y, \beta \rightarrow (\forall u, \exists v, \gamma))]$}
\begin{choices}
  \choice $[\exists x, \alpha \rightarrow (\forall y, \beta \rightarrow (\exists u, \forall v, \gamma))]$
  \choice $[\exists x, \alpha \rightarrow (\forall y, \beta \rightarrow (\exists u, \forall v, \neg \gamma))]$
  \choice $[\forall x, \neg \alpha \rightarrow (\exists y, \neg \beta \rightarrow (\forall u, \exists v, \neg \gamma))]$
  \choice $\left[ \exists x,\ \alpha \land \left( \forall y,\ \beta \land \left( \exists u,\ \forall v,\ \neg \gamma \right) \right) \right]$
\end{choices}

\end{multicols*}
\vspace{0.6cm}
\newpage
\begin{multicols*}{2}

\question {\small \textbf{[GATE 2008 | 2 Mark]}} \\[0.2cm]
{\Large P and Q are two propositions. which of the following logical expressions are equivalent?\newline
1. $P \vee \sim Q$\newline
2. $\sim (\sim P \wedge Q)$\newline
3. $(P \wedge Q) \vee (P \wedge \sim Q) \vee (\sim P \wedge \sim Q)$\newline
4. $(P \wedge Q) \vee (P \wedge \sim Q) \vee (\sim P \wedge Q)$}
\begin{choices}
  \choice Only 1 and 2
  \choice Only 1, 2 and 3
  \choice Only 1, 2 and 4
  \choice All of these
\end{choices}

\end{multicols*}
\vspace{0.6cm}
\newpage
\begin{multicols*}{2}

\question {\small \textbf{[GATE 2009 | 2 Mark]}} \\[0.2cm]
{\Large Which one of the following is the most appropriate logical formula to represent the statement: "Gold and silver ornaments are precious"\newline
The following notations are used:\newline
$G(x)$: $x$ is a gold ornament\newline
$S(x)$: $x$ is a silver ornament\newline
$P(x)$: $x$ is precious}
\begin{choices}
  \choice $\forall x(P(x) \rightarrow (G(x) \wedge S(x)))$
  \choice $\forall x(G(x) \wedge (S(x) \rightarrow P(x)))$
  \choice $\exists x((G(x) \wedge S(x)) \rightarrow P(x))$
  \choice $\forall x((G(x) \vee S(x)) \rightarrow P(x))$
\end{choices}

\end{multicols*}
\vspace{0.6cm}
\newpage
\begin{multicols*}{2}

\question {\small \textbf{[GATE 2009 | 2 Mark]}} \\[0.2cm]
{\Large The binary operation $\Box$ is defined as follows:\newline

\begin{center}
\includegraphics[max width=0.9\linewidth, keepaspectratio]{"E:/gate-pyq-app/public/1. Discrete Maths/1. Mathematical Logic/29.png"}
\end{center}
\newline
Which one of the following is equivalent to $P \vee Q$?}
\begin{choices}
  \choice $\neg Q \Box \neg P$
  \choice $P \Box \neg Q$
  \choice $\neg P \Box Q$
  \choice $\neg P \Box \neg Q$
\end{choices}

\end{multicols*}
\vspace{0.6cm}
\newpage
\begin{multicols*}{2}

\question {\small \textbf{[GATE 2009 | 2 Mark]}} \\[0.2cm]
{\Large Consider the following well-formed formulae:\newline
I. $\neg \forall x(P(x))$\newline
II. $\neg \exists x(P(x))$\newline
III. $\neg \exists x(\neg P(x))$\newline
IV. $\exists x(\neg P(x))$\newline
Which of the above are equivalent?}
\begin{choices}
  \choice I and III
  \choice I and IV
  \choice II and III
  \choice II and IV
\end{choices}

\end{multicols*}
\vspace{0.6cm}
\newpage
\begin{multicols*}{2}

\question {\small \textbf{[GATE 2010 | 2 Mark]}} \\[0.2cm]
{\Large Suppose the predicate $F(x, y, t)$ is used to represent the statement that person $x$ can fool person $y$ at time $t$. Which one of the statements below expresses best the meaning of the formula $\forall x \exists y \exists t (\neg F(x, y, t))$?}
\begin{choices}
  \choice Everyone can fool some person at some time
  \choice No one can fool everyone all the time
  \choice Everyone cannot fool some person all the time
  \choice No one can fool some person at some time
\end{choices}

\end{multicols*}
\vspace{0.6cm}
\newpage
\begin{multicols*}{2}

\question {\small \textbf{[GATE 2011 | 2 Mark]}} \\[0.2cm]
{\Large Which one of the following options is CORRECT given three positive integers $x, y$ and $z$, and a predicate\newline
$P(x) = \neg (x = 1) \wedge \forall y (\exists z (x = y * z) \Rightarrow (y = x) \vee (y = 1))$}
\begin{choices}
  \choice $P(x)$ being true means that $x$ is a prime number
  \choice $P(x)$ being true means that $x$ is a number other than 1
  \choice $P(x)$ is always true irrespective of the value of $x$
  \choice $P(x)$ being true means that $x$ has exactly two factors other than 1 and $x$
\end{choices}

\end{multicols*}
\vspace{0.6cm}
\newpage
\begin{multicols*}{2}

\question {\small \textbf{[GATE 2012 | 1 Mark]}} \\[0.2cm]
{\Large Consider the following logical inferences.\newline
$I_1$: If it rains then the cricket match will not be played.\newline
The cricket match was played.\newline
Inference: There was no rain.\newline
$I_2$: If it rains then the cricket match will not be played.\newline
It did not rain.\newline
Inference: The cricket match was played.\newline
Which of the following is TRUE?}
\begin{choices}
  \choice Both $I_1$ and $I_2$ are correct inferences
  \choice $I_1$ is correct but $I_2$ is not a correct inference
  \choice $I_1$ is not correct but $I_2$ is a correct inference
  \choice Both $I_1$ and $I_2$ are not correct inferences
\end{choices}

\end{multicols*}
\vspace{0.6cm}
\newpage
\begin{multicols*}{2}

\question {\small \textbf{[GATE 2012 | 1 Mark]}} \\[0.2cm]
{\Large What is the correct translation of the following statement into mathematical logic?\newline
"Some real numbers are rational"}
\begin{choices}
  \choice $\exists x (\text{real} (x) \vee \text{rational} (x))$
  \choice $\forall x (\text{real} (x) \rightarrow \text{rational} (x))$
  \choice $\exists x (\text{real} (x) \wedge \text{rational} (x))$
  \choice $\exists x (\text{rational} (x) \rightarrow \text{real} (x))$
\end{choices}

\end{multicols*}
\vspace{0.6cm}
\newpage
\begin{multicols*}{2}

\question {\small \textbf{[GATE 2013 | 2 Mark]}} \\[0.2cm]
{\Large What is the logical translation of the following statements?\newline
"None of my friends are perfect"}
\begin{choices}
  \choice $\exists x (F(x) \wedge \neg P(x))$
  \choice $\exists x (\neg F(x) \wedge P(x))$
  \choice $\exists x (\neg F(x) \wedge \neg P(x))$
  \choice $\neg \exists x (F(x) \wedge P(x))$
\end{choices}

\end{multicols*}
\vspace{0.6cm}
\newpage
\begin{multicols*}{2}

\question {\small \textbf{[GATE 2013 | 2 Mark]}} \\[0.2cm]
{\Large Which one of the following is NOT logically equivalent to $\neg \exists x(\forall y(\alpha) \wedge \forall z(\beta))$}
\begin{choices}
  \choice $\forall x(\exists z(\neg \beta) \rightarrow \forall y(\alpha))$
  \choice $\forall x(\forall z(\beta) \rightarrow \exists y(\neg \alpha))$
  \choice $\forall x(\forall y(\alpha) \rightarrow \exists z(\neg \beta))$
  \choice $\forall x(\exists y(\neg \alpha) \vee \exists z(\neg \beta))$
\end{choices}

\end{multicols*}
\vspace{0.6cm}
\newpage
\begin{multicols*}{2}

\question {\small \textbf{[GATE 2014 | 1 Mark]}} \\[0.2cm]
{\Large Consider the statement:\newline
"Not all that glitters is gold"\newline
Predicate $glitters(x)$ is true if $x$ glitters and predicate $gold(x)$ is true if $x$ is gold. Which one of the following logical formulae represents the above statement?}
\begin{choices}
  \choice $\forall x: glitters(x) \Rightarrow \neg gold(x)$
  \choice $\forall x: gold(x) \Rightarrow glitters(x)$
  \choice $\exists x: gold(x) \wedge \neg glitters(x)$
  \choice $\exists x: glitters(x) \wedge \neg gold(x)$
\end{choices}

\end{multicols*}
\vspace{0.6cm}
\newpage
\begin{multicols*}{2}

\question {\small \textbf{[GATE 2014 | 2 Mark]}} \\[0.2cm]
{\Large Which one of the following propositional logic formulas is TRUE when exactly two of $p, q$, and $r$ are TRUE?}
\begin{choices}
  \choice $((p \leftrightarrow q) \wedge r) \vee (p \wedge q \wedge \sim r)$
  \choice $(\sim(p \leftrightarrow q) \wedge r) \vee (p \wedge q \wedge \sim r)$
  \choice $((p \rightarrow q) \wedge r) \vee (p \wedge q \wedge \sim r)$
  \choice $(\sim(p \leftrightarrow q) \wedge r) \wedge (p \wedge q \wedge \sim r)$
\end{choices}

\end{multicols*}
\vspace{0.6cm}
\newpage
\begin{multicols*}{2}

\question {\small \textbf{[GATE 2014 | 2 Mark]}} \\[0.2cm]
{\Large Which one of the following Boolean expressions is NOT a tautology?}
\begin{choices}
  \choice $((a \rightarrow b) \wedge (b \rightarrow c)) \rightarrow (a \rightarrow c)$
  \choice $(a \leftrightarrow c) \rightarrow (\sim b \rightarrow (a \wedge c))$
  \choice $(a \wedge b \wedge c) \rightarrow (c \vee a)$
  \choice $a \rightarrow (b \rightarrow a)$
\end{choices}

\end{multicols*}
\vspace{0.6cm}
\newpage
\begin{multicols*}{2}

\question {\small \textbf{[GATE 2014 | 1 Mark]}} \\[0.2cm]
{\Large Consider the following statements:\newline
P: Good mobile phones are not cheap\newline
Q: Cheap mobile phones are not good\newline
L: P implies Q\newline
M: Q implies P\newline
N: P is equivalent to Q\newline
Which one of the following about L, M, and N is CORRECT?}
\begin{choices}
  \choice Only L is TRUE
  \choice Only M is TRUE
  \choice Only N is TRUE
  \choice L, M and N are TRUE
\end{choices}

\end{multicols*}
\vspace{0.6cm}
\newpage
\begin{multicols*}{2}

\question {\small \textbf{[GATE 2014 | 2 Mark]}} \\[0.2cm]
{\Large The CORRECT formula for the sentence, "not all rainy days are cold" is}
\begin{choices}
  \choice $\forall d (\text{Rainy}(d) \wedge \sim \text{Cold}(d))$
  \choice $\forall d (\sim \text{Rainy}(d) \rightarrow \text{Cold}(d))$
  \choice $\exists d (\sim \text{Rainy}(d) \rightarrow \text{Cold}(d))$
  \choice $\exists d (\text{Rainy}(d) \wedge \sim \text{Cold}(d))$
\end{choices}

\end{multicols*}
\vspace{0.6cm}
\newpage
\begin{multicols*}{2}

\question {\small \textbf{[GATE 2015 | 1 Mark]}} \\[0.2cm]
{\Large Which one of the following is Not equivalent to $p \leftrightarrow q$?}
\begin{choices}
  \choice $(\sim p \vee q) \wedge (p \vee \sim q)$
  \choice $(\sim p \vee q) \wedge (q \rightarrow p)$
  \choice $(\sim p \wedge q) \vee (p \wedge \sim q)$
  \choice $(\sim p \wedge \sim q) \vee (p \wedge q)$
\end{choices}

\end{multicols*}
\vspace{0.6cm}
\newpage
\begin{multicols*}{2}

\question {\small \textbf{[GATE 2015 | 2 Mark]}} \\[0.2cm]
{\Large The binary operator $\oplus$ is defined by the following truth table.\newline

\begin{center}
\includegraphics[max width=0.9\linewidth, keepaspectratio]{"E:/gate-pyq-app/public/1. Discrete Maths/1. Mathematical Logic/43.png"}
\end{center}
\newline
Which one of the following is true about the binary operator $\oplus$?}
\begin{choices}
  \choice Both commutative and associative
  \choice Commutative but not associative
  \choice Not commutative but associative
  \choice Neither commutative nor associative
\end{choices}

\end{multicols*}
\vspace{0.6cm}
\newpage
\begin{multicols*}{2}

\question {\small \textbf{[GATE 2015 | 1 Mark]}} \\[0.2cm]
{\Large Consider the following two statements:\newline
$S_1$: If a candidate is known to be corrupt, then he will not be elected.\newline
$S_2$: If a candidate is kind, he will be elected.\newline
Which one of the following statements follows from $S_1$ and $S_2$ as per sound inference rules of logic?}
\begin{choices}
  \choice If a person is known to be corrupt, he is kind
  \choice If a person is not known to be corrupt, he is not kind
  \choice If a person is kind, he is not known to be corrupt
  \choice If a person is not kind, he is not known to be corrupt
\end{choices}

\end{multicols*}
\vspace{0.6cm}
\newpage
\begin{multicols*}{2}

\question {\small \textbf{[GATE 2015 | 2 Mark]}} \\[0.2cm]
{\Large Which one of the following well formed formulae is a tautology?}
\begin{choices}
  \choice $\forall x \exists y R(x,y) \leftrightarrow \exists y \forall x R(x,y)$
  \choice $(\forall x [\exists y R(x,y) \rightarrow S(x,y)]) \rightarrow \forall x \exists y S(x,y)$
  \choice $[\forall x \exists y (P(x,y) \rightarrow R(x,y))] \leftrightarrow [\forall x \exists y (\neg P(x,y) \vee R(x,y))]$
  \choice $\forall x \forall y P(x,y) \rightarrow \forall x \forall y P(y,x)$
\end{choices}

\end{multicols*}
\vspace{0.6cm}
\newpage
\begin{multicols*}{2}

\question {\small \textbf{[GATE 2015 | 1 Mark]}} \\[0.2cm]
{\Large In a room there are only two types of people, namely Type 1 and Type 2. Type 1 people always tell the truth and Type 2 people always lie. You give a fair coin to a person in that room, without knowing which type he is from and tell him to toss it and hide the result from you till you ask for it. Upon asking, the person replies the following:\newline
$\textit{``The result of the toss is head if and only if I am telling the truth.``}$\newline
Which of the following options is correct?}
\begin{choices}
  \choice The result is head
  \choice The result is tail
  \choice If the person is of Type 2, then the result is tail
  \choice If the person is of Type 1, then the result is tail
\end{choices}

\end{multicols*}
\vspace{0.6cm}
\newpage
\begin{multicols*}{2}

\question {\small \textbf{[GATE 2016 | 1 Mark]}} \\[0.2cm]
{\Large Let $p, q, r, s$ represent the following propositions.\newline
$p$: $x \in \{8, 9, 10, 11, 12\}$\newline
$q$: $x$ is a composite number\newline
$r$: $x$ is a perfect square\newline
$s$: $x$ is a prime number\newline
The integer $x \geq 2$ which satisfies $\neg ((p \Rightarrow q) \wedge (\sim r \vee \sim s))$ is \rule{2cm}{0.4pt}}

\vspace{0.5cm}
\textbf{Answer:} \rule{3cm}{0.4pt}

\end{multicols*}
\vspace{0.6cm}
\newpage
\begin{multicols*}{2}

\question {\small \textbf{[GATE 2016 | 1 Mark]}} \\[0.2cm]
{\Large Consider the following expressions:\newline
(i) False\newline
(ii) $Q$\newline
(iii) True\newline
(iv) $P \vee Q$\newline
(v) $\neg Q \vee P$\newline
The number of expressions given above that are logically implied by $P \wedge (P \Rightarrow Q)$ is \rule{2cm}{0.4pt}}

\vspace{0.5cm}
\textbf{Answer:} \rule{3cm}{0.4pt}

\end{multicols*}
\vspace{0.6cm}
\newpage
\begin{multicols*}{2}

\question {\small \textbf{[GATE 2016 | 2 Mark]}} \\[0.2cm]
{\Large Which one of the following well-formed formulae in predicate calculus is NOT valid?}
\begin{choices}
  \choice $(\forall x p(x) \Rightarrow \forall x q(x)) \Rightarrow (\exists x \neg p(x) \vee \forall x q(x))$
  \choice $(\exists x p(x) \vee \exists x q(x)) \Rightarrow \exists x (p(x) \vee q(x))$
  \choice $(\exists x (p(x) \wedge q(x))) \Rightarrow (\exists x p(x) \wedge \exists x q(x))$
  \choice $\forall x (p(x) \vee q(x)) \Rightarrow (\forall x p(x) \vee \forall x q(x))$
\end{choices}

\end{multicols*}
\vspace{0.6cm}
\newpage
\begin{multicols*}{2}

\question {\small \textbf{[GATE 2017 | 1 Mark]}} \\[0.2cm]
{\Large Consider the first-order logic sentence $\forall x \exists y R(x,y)$. Assuming non-empty logical domains, which of the sentences below are implied by $F$?\newline
I. $\exists y \forall x R(x,y)$\newline
II. $\exists y \forall x R(x,y)$\newline
III. $\forall y \exists x R(x,y)$\newline
IV. $\exists x \forall y \neg R(x,y)$}
\begin{choices}
  \choice IV only
  \choice I and IV only
  \choice II only
  \choice II and III only
\end{choices}

\end{multicols*}
\vspace{0.6cm}
\newpage
\begin{multicols*}{2}

\question {\small \textbf{[GATE 2017 | 1 Mark]}} \\[0.2cm]
{\Large The statement $(\neg p) \Rightarrow (\neg q)$ is logically equivalent to which of the statements below?\newline
I. $p \Rightarrow q$\newline
II. $q \Rightarrow p$\newline
III. $(\neg q) \vee p$\newline
IV. $(\neg p) \vee q$}
\begin{choices}
  \choice I only
  \choice I and IV only
  \choice II only
  \choice II and III only
\end{choices}

\end{multicols*}
\vspace{0.6cm}
\newpage
\begin{multicols*}{2}

\question {\small \textbf{[GATE 2026 | 1 Mark]}} \\[0.2cm]
{\Large Let $p, q,$ and $r$ be propositions and the expression $(p \Rightarrow q) \Rightarrow r$ be a contradiction. Then, the expression $(r \Rightarrow p) \Rightarrow \neg q$ is}
\begin{choices}
  \choice A tautology
  \choice A contradiction
  \choice Always TRUE when $p$ is FALSE
  \choice Always TRUE when $q$ is TRUE
\end{choices}

\end{multicols*}
\vspace{0.6cm}
\newpage
\begin{multicols*}{2}

\question {\small \textbf{[GATE 2026 | 1 Mark]}} \\[0.2cm]
{\Large Let $p, q, r$ denote the statements “It is raining”, “It is cold”, and “It is pleasant”, respectively. Then the statement “It is not raining and it is pleasant, and it is not pleasant only if it is raining and it is cold” is represented by}
\begin{choices}
  \choice $(\neg p \wedge r) \wedge (\neg r \Rightarrow (p \wedge q))$
  \choice $(\neg p \wedge r) \wedge ((p \wedge q) \Rightarrow \neg r)$
  \choice $(\neg p \wedge r) \vee ((p \wedge q) \Rightarrow \neg r)$
  \choice $(\neg p \wedge r) \vee (\neg r \Rightarrow (p \wedge q))$
\end{choices}

\end{multicols*}
\vspace{0.6cm}
\newpage
\begin{multicols*}{2}

\question {\small \textbf{[GATE 2019 | 2 Mark]}} \\[0.2cm]
{\Large Consider this first order predicate formula\newline
$\forall x[(\forall z; z \mid x \Rightarrow (z = x \vee z = 1)) \Rightarrow \exists u (w > x) \vee (w = z) \vee (z = 1)]$\newline
Here $a \mid b$ denotes that “$a$ divides $b$”, where $a$ and $b$ are integers. Consider the following sets:\newline
$S_1$: ${1, 2, 3, ..., 100}$\newline
$S_2$: Set of all positive integers\newline
$S_3$: Set of all integers\newline
Which of the above sets satisfy $\varphi$?}
\begin{choices}
  \choice $S_1$ and $S_3$
  \choice $S_2$ and $S_3$
  \choice $S_1, S_2$ and $S_3$
  \choice $S_1$ and $S_2$
\end{choices}

\end{multicols*}
\vspace{0.6cm}
\newpage
\begin{multicols*}{2}

\question {\small \textbf{[GATE 2020 | 2 Mark]}} \\[0.2cm]
{\Large Which one of the following predicate formulae is NOT logically valid? Note that $W$ is a predicate formula without any free occurrence of $x$.}
\begin{choices}
  \choice $\exists x (p(x) \wedge W) \equiv \exists x p(x) \wedge W$
  \choice $\forall x (p(x) \Rightarrow W) \equiv \forall x p(x) \Rightarrow W$
  \choice $\exists x (p(x) \Rightarrow W) \equiv \forall x p(x) \Rightarrow W$
  \choice $\forall x (p(x) \vee W) \equiv \forall x p(x) \vee W$
\end{choices}

\end{multicols*}
\vspace{0.6cm}
\newpage
\begin{multicols*}{2}

\question {\small \textbf{[GATE 2021 | 1 Mark]}} \\[0.2cm]
{\Large Let $p$ and $q$ be two propositions. Consider the following formulae in propositional logic.\newline
$S_1: (\neg p \vee q) \Rightarrow q$\newline
$S_2: q \Rightarrow (p \vee q)$\newline
Which one of the following choices is correct?}
\begin{choices}
  \choice Neither $S_1$ nor $S_2$ is a tautology
  \choice Both $S_1$ and $S_2$ are tautologies
  \choice $S_1$ is a tautology but $S_2$ is not a tautology
  \choice $S_1$ is not a tautology but $S_2$ is a tautology
\end{choices}

\end{multicols*}
\vspace{0.6cm}
\newpage
\begin{multicols*}{2}

\question {\small \textbf{[GATE 2021 | 1 Mark]}} \\[0.2cm]
{\Large Choose the correct choice(s) regarding the following propositional logic assertion $S$:\newline
$S: ((P \wedge Q) \Rightarrow R) \Rightarrow ((P \wedge Q) \Rightarrow (Q \Rightarrow R))$}
\begin{choices}
  \choice $S$ is a tautology
  \choice $S$ is neither a tautology nor a contradiction
  \choice The antecedent of $S$ is logically equivalent to consequent of $S$
  \choice $S$ is a contradiction
\end{choices}

\end{multicols*}
\vspace{0.6cm}
\newpage
\begin{multicols*}{2}

\question {\small \textbf{[GATE 2023 | 1 Mark]}} \\[0.2cm]
{\Large Geetha has a conjecture about integers, which is of the form $\forall x (P(x) \Rightarrow \exists y Q(x,y))$, where $P$ is a statement about integers, and $Q$ is a statement about pairs of integers. Which of the following (one or more) option(s) imply Geetha’s conjecture?}
\begin{choices}
  \choice $\exists x (P(x) \wedge \forall y Q(x,y))$
  \choice $\forall x \forall y Q(x,y)$
  \choice $\exists y \forall x (P(x) \Rightarrow Q(x,y))$
  \choice $\exists x (P(x) \wedge \exists y Q(x,y))$
\end{choices}

\end{multicols*}
\vspace{0.6cm}
\newpage
\begin{multicols*}{2}

\question {\small \textbf{[GATE 2024 | 1 Mark]}} \\[0.2cm]
{\Large Let $p$ and $q$ be the following propositions:\newline
$p$: Fail grade can be given.\newline
$q$: Student scores more than $50\%$ marks.\newline
Consider the statement “Fail grade cannot be given when student scores more than $50\%$ marks.” Which one of the following is the CORRECT representation of the above statement in propositional logic?}
\begin{choices}
  \choice $q \Rightarrow p$
  \choice $q \Rightarrow \neg p$
  \choice $p \Rightarrow q$
  \choice $\neg p \Rightarrow \neg q$
\end{choices}

\end{multicols*}
\vspace{0.6cm}
\newpage

\end{questions}
\end{document}
